\section{Matched Ordered-CDF Output Geometry}
\label{app:ordered-cdf-geometry}

Consider probability vectors on an ordered grid with $m$ bins. Let
\[
P=I-\frac{1}{m}\mathbf 1\mathbf 1^\top
\]
project onto the probability-simplex tangent, let
$C\in\mathbb R^{(m-1)\times m}$ compute cumulative sums through the
penultimate bin, and let $\Delta x$ be the grid spacing. Define
\[
G_{\mathrm{CDF}}=\Delta x\,P C^\top C P.
\]
For a zero-sum probability perturbation $\Delta z$,
\[
\Delta z^\top G_{\mathrm{CDF}}\Delta z
=
\Delta x\,\|C\Delta z\|_2^2,
\]
which is the squared $L^2$ distance between the corresponding
one-dimensional cumulative distribution functions. Because
$G_{\mathrm{CDF}}$ is singular in the ambient space, its use is restricted to
the probability-simplex tangent. The matched dual form for a loss covector
$r$ is
\[
(Pr)^\top G_{\mathrm{CDF}}^+(Pr),
\]
where $G_{\mathrm{CDF}}^+$ is the Moore--Penrose pseudoinverse. The
implementation tests both the primal identity and the matched dual
construction numerically.

This is an ordered-CDF or Cram\'er metric, not sliced Wasserstein. It supplies
one explicit space-dependent geometry compatible with IEWC's matched
primal/dual construction. The retained diffusion output-representation
experiment predates this correction and is therefore treated only as
exploratory evidence that alternative output representations can affect the
penalty, not as empirical validation of Wasserstein geometry.
